\documentclass[11pt]{amsart}
\usepackage[margin=1.15in]{geometry}
\usepackage{amsmath,amssymb,amsthm,booktabs}
\usepackage[colorlinks=true,allcolors=blue]{hyperref}
\newtheorem{theorem}{Theorem}
\newtheorem{proposition}[theorem]{Proposition}
\theoremstyle{remark}
\newtheorem{remark}[theorem]{Remark}
\begin{document}
\title[Three remarks on Pre-Fascicle 8A]{Three remarks on TAOCP Pre-Fascicle 8A: exercises 157(b), 204(b) and 205}
\author{Zijian Zeng}
\address{UCSI University, Kuala Lumpur, Malaysia}
\date{18 November 2026}
\begin{abstract}
We explain why no closed $8\times8$ knight's tour has at least two moves in each of the $21$ classes of Knuth's
exercise 7.2.2.4--157(b), for which the draft answer says that no explanation is known: five of the classes
together always contain exactly eight moves. We point out that the field count in the trie-building loop of
answer 204(b) has its two cases interchanged. And we solve the case $n=1$ of the research problem 205, which the
answer suggests should be solved first: the exact average number of pointers for one random key, with its
asymptotic expansion.
\end{abstract}
\maketitle

All references are to Knuth's draft of Section 7.2.2.4 (Hamiltonian paths and cycles), Pre-Fascicle 8A of
\emph{The Art of Computer Programming}, dated 8 May 2026~\cite{knuth8a}; exercise and answer numbers are those of
that draft.

\section{Exercise 157(b)}
Exercise 157 partitions the moves of the $8\times8$ knight graph into $21$ classes, the orbits under the eight
symmetries of the board, named by letters in a figure; part (b) asks for the closed tours that use at least two
moves of every class, and the answer is that there are none, with no explanation known. Name the cells $ij$
($0\le i,j<8$, row then column). Reading the letters from the figure, the six classes that matter here are those
of the moves $A=00\text{--}12$, $D=03\text{--}15$, $I=12\text{--}24$, $K=12\text{--}31$, $L=02\text{--}21$,
$N=12\text{--}33$.

\begin{theorem}
In every closed knight's tour of the $8\times8$ board, the classes $D$, $I$, $K$, $L$, $N$ together contain exactly
eight moves. Hence one of them contains at most one move, and no tour has two moves of every class.
\end{theorem}
\begin{proof}
Let $C$ be the eight cells a knight's move from a corner, namely $12,21,15,26,51,62,56,65$. A corner has only two
moves, so a closed tour uses both of them, and each cell of $C$ spends one of its two tour moves on its corner (a
move of class $A$). Its other tour move goes to another neighbour. The moves at cell $12$ other than $12\text{--}00$
are $12\text{--}04$, $12\text{--}20$, $12\text{--}24$, $12\text{--}31$, $12\text{--}33$, of classes $D$, $L$, $I$, $K$,
$N$; by symmetry the same holds at every cell of $C$. Finally, every move of these five classes has exactly one end
in $C$. So each cell of $C$ contributes exactly one move of the five classes, and no move is counted twice.
\end{proof}

The same argument works on every $n\times n$ board with $n\ge7$ (checked for $7\le n\le20$; for larger $n$ the
neighbourhoods of the corners are far apart). It fails for $n=6$, where two cells next to different corners are a
knight's move apart. A CP-SAT computation, removing classes one at a time from all $21$, arrives at
$\{D,I,K,L,N\}$ as a minimal set of classes whose ``at least two moves each'' requirement is already infeasible;
in $20$ random closed tours the five classes carried exactly eight moves.

\section{Answer 204(b)}
Answer 204(b) builds compact tries whose nodes have a variable number of pointer fields. After reading $a_l$ there are
$\mathrm{TMS}$ open pairs. If $l+\mathrm{TMS}=q$ (tight), only the $\mathrm{TMS}$ closing fields are needed. Otherwise
the node has fields for $\bar1$, $0$, the $\mathrm{TMS}$ closings, and one for opening a new pair; but a new pair can be
opened only if $q-l-\mathrm{TMS}\ge2$. So the count should be
\[d\gets\mathrm{TMS}+(l+\mathrm{TMS}=q-1\ ?\ 2:3),\]
whereas the draft has $(\dots ?\ 3:2)$. This agrees with the prose of the same answer, which says that for the prefix
$1201\bar132$ the field `4' is unnecessary when $q=8$ or $9$. Running the loop literally as printed on $300$ random
keys, $14$, $86$, $201$ and $484$ field accesses fall outside the node just allocated for $q=6,8,10,14$; with the two
cases interchanged there are none.

\section{Exercise 205 with one key}
Exercise 205 asks for the average total number of pointers when the method of exercise 204 builds a compact trie
from $n$ random keys; each key is a uniformly random marked involution of order $q$, written $a_1\ldots a_q$ with
$\bar1$ and $0$ for the two kinds of fixed points and codes $1,2,\dots$ for the pairs in order of their first
element. The answer suggests solving the case $n=1$ first, with its asymptotics as $q\to\infty$. Let $T_q$ be the
number of marked involutions ($T_q=2T_{q-1}+(q-1)T_{q-2}$).

With one key the trie has a root of three fields and one node for each prefix of length $1\le l<q$, with
$d_l=\mathrm{TMS}_l+3$ fields, minus $1$ if $l+\mathrm{TMS}_l=q-1$ and minus $3$ if $l+\mathrm{TMS}_l=q$ (using the
corrected rule of \S2).

\begin{proposition}
The average number of pointers for one random key of length $q\ge2$ is
\[E_q=3q+\frac{(q+1)(q-\bar F)}{6}-\sum_{s=1}^{q-1}\frac{1}{T_q}\Bigl(\frac{2s\,(q-s)!}{(q-2s+1)!}\,T_{q-2s+1}
+\frac{3\,(q-s)!}{(q-2s)!}\,T_{q-2s}\Bigr),\qquad \bar F=\frac{2qT_{q-1}}{T_q},\]
where terms with a negative factorial are omitted. For example $E_2=24/5$, $E_3=54/7$, $E_4=471/43$.
\end{proposition}
\begin{proof}
$\mathrm{TMS}_l$ is the number of pairs with one element in $\{1,\dots,l\}$ and the other after it. A pair at
positions $i<j$ is counted for $j-i$ values of $l$, and $E(j-i)=(q+1)/3$ for a random pair; the expected number of
pairs is $(q-\bar F)/2$, where $\bar F=q\cdot2T_{q-1}/T_q$ is the expected number of fixed points. With $s=q-l$,
the case $l+\mathrm{TMS}_l=q$ means that all $s$ elements after position $l$ are second elements of pairs whose
first element is at most $l$; there are $(q-s)!/(q-2s)!$ ways to choose those partners and $T_{q-2s}$ ways to finish.
The case $l+\mathrm{TMS}_l=q-1$ means the same except that exactly one of the $s$ elements is a fixed point ($2s$
choices). Summing gives the formula.
\end{proof}
The formula agrees with the average over all $T_q$ keys for $2\le q\le9$.

\begin{proposition}
As $q\to\infty$,
\[E_q=\frac{q^2}{6}-\frac{q^{3/2}}{3}+\frac{7q}{2}-c\sqrt q+O(1),\qquad
c=\frac{31}{12}-e^2\sqrt{\pi/2}\,\operatorname{erfc}\sqrt2\approx2.1619.\]
\end{proposition}
\noindent\emph{Derivation.} From the recurrence, $T_q/T_{q-1}=\sqrt q+1+1/(4\sqrt q)+O(1/q)$, hence
$\bar F=2\sqrt q-2+\frac32q^{-1/2}+O(1/q)$, which gives the first three terms and $-\frac7{12}\sqrt q$. In the sums,
with $s=x\sqrt q$ the tight term is $e^{-x^2/2-2x}(1+o(1))$ and the near-tight term is $2x$ times that, so they
contribute $-\sqrt q\bigl(3I+2(1-2I)\bigr)=-(2-I)\sqrt q$, where $I=\int_0^\infty e^{-x^2/2-2x}\,dx=
e^2\sqrt{\pi/2}\,\operatorname{erfc}\sqrt2\approx0.42137$. The first three terms follow rigorously from the exact
formula; the constant $c$ comes from this standard Laplace estimate, and the exact values give
$(E_q-q^2/6+q^{3/2}/3-7q/2)/\sqrt q=-2.073,\,-2.117,\,-2.139$ for $q=400,1600,6400$, approaching $-c$.

The case of $n\ge2$ keys, where prefixes are shared, is not treated here.

\section*{Checks}
The program \texttt{notes8a\_check.py} (Python standard library only, a few seconds) checks the two structural facts
of \S1 for $7\le n\le20$, runs the loop of \S2 as printed and corrected, and compares the formula of \S3 with the
brute-force average for $q\le9$. The CP-SAT computations of \S1 are in \texttt{search/}.

\medskip\noindent\emph{How this was done.} The computations, the arguments and this note were produced with Claude
Code (Anthropic) at the author's direction. The author is responsible for the content.

\begin{thebibliography}{9}
\bibitem{knuth8a} D.~E. Knuth, \emph{The Art of Computer Programming}, Pre-Fascicle 8A: Hamiltonian paths and cycles
(Section 7.2.2.4), draft of 8 May 2026.
\end{thebibliography}
\end{document}
